% !TeX root = ./settheory-screen.tex
% settheory.tex
%
% driver file settheory.tex to produce text

\preto\OLEndChapterHook{\IfFileExists{include/summary-\thechapter}
        {\section*{小结}\addcontentsline{toc}{section}{小结}
        \let\emph\textbf\input{include/summary-\thechapter}\let\emph\textit}{}}

\problemsperchapter

\allowdisplaybreaks

\frontmatter

\OLPfrontmatter

\chapter{关于本书}

本书是一份开放教育资源。它面向稍有逻辑学背景和高中数学知识的学生。其目标是浅尝集合论哲学的皮毛。读完本书，学生可能会对以下几点有所认识：

\begin{enumerate}
	\item 集合论是如何产生的；
	\item 如何在集合论加算术的框架下嵌入数学的大部分领域；
	\item 如何将算术本身嵌入集合论；
	\item 集合的累积迭代概念意味着什么；
	\item 如何尝试为 ZFC 的公理辩护。
\end{enumerate}

本书脱胎于我在剑桥大学哲学系讲授的一门短期课程。在我之前，这门课由 Luca Incurvati（卢卡·因库尔瓦蒂）和 Michael Potter（迈克尔·波特）主讲。在撰写本书——以及更广泛地，准备这门课程——的过程中，我深感受惠于 Luca 和 Michael 两人。我希望这份感激之情能在本书的正文中体现出来；但在此，我也要向他们表达由衷的感谢。

本书的大部分内容最初以《Open Set Theory》之名发布；本书是其续作。我已将本书所有材料贡献给 \href{http://openlogicproject.org}{\emph{Open Logic Project}}，在那里可以免费获取。（\Crefrange{sfr:set::chap}{sfr:siz::chap} 取材自 OLP 中已有的材料，仅作了细微改动；这些改动现在也已成为 OLP 的一部分。）欲了解更多信息，请访问 \href{http://openlogicproject.org/}{openlogicproject.org}。

\begin{flushright}
Tim Button（蒂姆·巴顿）
\\伦敦大学学院
\\2021年10月
\end{flushright}

\mainmatter

\part{序篇}

\olchapter{his}{set}{历史与神话}

要理解集合论的哲学意义，有必要先了解集合论究竟为何会产生。而要了解这一点，不妨稍许思考数学的历史与神话。因此，在开始讨论集合论本身之前，我们将从一个非常简短的“历史”讲起。但我们给这个词加上了引号，因为它非常简略，极具选择性，并且多少有些争议。

\olimport*[history/set-theory]{infinitesimals}
\olimport*[history/set-theory]{limits}
\olimport*[history/set-theory]{pathologies}
\olimport*[history/set-theory]{mythology}

\section{路线图}

上一节的部分寓意在于，数学史在很大程度上是由胜利者书写的。他们各有自己的主张要推行——哲学的和数学的主张。严肃地研究数学史极其困难（也回报丰厚），而密涅瓦的猫头鹰只在黄昏时分起飞。

尽管如此，无可争议的是，那些“病态的”结果推动了令人着迷的新数学工具的发展，并促使人们重新思考数学严谨性的标准。例如，它们要求以特定的方式思考连续统（“实数线”），并将函数视为逐点的映射。而最终，所有这些工具的全面发展，都需要严格地发展集合论。本书余下部分将介绍这一发展的部分内容。

\Olref[sfr][][]{part}将介绍一种\emph{朴素}集合论，它足以发展出刚才描述的全部数学。这部分内容将占用不少篇幅。但到\olref[sfr][][]{part}结束时，我们将能够理解如何将实数视为某些集合，以及如何将实数上的函数（包括充满空间的曲线）视为\emph{更进一步}的集合。

但这种集合论的\emph{朴素性}将在\olref[sth][][]{part}中显现出来，届时我们将遇到集合论悖论，并感到需要更精确地描述事物。此时，我们将需要发展一套关于集合的公理化处理方式，用以重获我们所有的朴素结果，同时（希望）避免悖论。（数学严谨性的猫头鹰，也只在黄昏时分起飞。）

\OLEndChapterHook

\olpart{sfr}{朴素集合论}

\section*{\Olref[sfr][][]{part}导言}

在\olref[sfr][][]{part}中，我们将以一种朴素、非形式化的方式考虑集合。\Olref[sfr][set][]{chap}将介绍基本思想，即集合是从外延角度考虑的汇集，并介绍一些非常基本的运算。接着，\crefrange{sfr:rel::chap}{sfr:fun::chap}将解释集合论如何让我们谈论关系以及（由此而来的）函数。

然后，\Crefrange{sfr:siz::chap}{sfr:infinite::chap}将讨论朴素集合论的一些早期成果。在\olref[sfr][siz][]{chap}中，我们探讨如何比较集合的大小。在\olref[sfr][arith][]{chap}中，我们探讨如何将整数、有理数和实数还原为集合论与基本算术。在\olref[sfr][infinite][]{chap}中，我们考虑如何在集合论内实现基本算术。

重申一下，所有这些工作都将以\emph{朴素}的方式进行。但我们在\olref[sfr][][]{part}中所做的一切，都可以在\olref[sth][][]{part}中引入的形式化集合论中，完全严格地实现。

\olchapter{sfr}{set}{入门}

\olimport*[sets-functions-relations/sets]{basics}
\olimport*[sets-functions-relations/sets]{subsets}
\olimport*[sets-functions-relations/sets]{important-sets}
\olimport*[sets-functions-relations/sets]{unions-and-intersections}
\olimport*[sets-functions-relations/sets]{pairs-and-products}
\olimport*[sets-functions-relations/sets]{russells-paradox}

\OLEndChapterHook

\olchapter{sfr}{rel}{关系}

\olimport*[sets-functions-relations/relations]{relations-as-sets}
\olimport*[sets-functions-relations/relations]{reflections}
\olimport*[sets-functions-relations/relations]{special-properties}
\olimport*[sets-functions-relations/relations]{equivalence-relations}
\olimport*[sets-functions-relations/relations]{orders}
\olimport*[sets-functions-relations/relations]{operations}

\OLEndChapterHook

\olchapter{sfr}{fun}{函数}

\olimport*[sets-functions-relations/functions]{function-basics}
\olimport*[sets-functions-relations/functions]{function-kinds}
\olimport*[sets-functions-relations/functions]{functions-relations}
\olimport*[sets-functions-relations/functions]{inverses}
\olimport*[sets-functions-relations/functions]{composition}

\OLEndChapterHook

\olchapter{sfr}{siz}{集合的大小}

\olimport*[sets-functions-relations/size-of-sets]{introduction}
\olimport*[sets-functions-relations/size-of-sets]{enumerability-alt}
\olimport*[sets-functions-relations/size-of-sets]{zig-zag}
\olimport*[sets-functions-relations/size-of-sets]{pairing}
\olimport*[sets-functions-relations/size-of-sets]{non-enumerability-alt}
\olimport*[sets-functions-relations/size-of-sets]{reduction-alt}
\olimport*[sets-functions-relations/size-of-sets]{equinumerous-sets}
\olimport*[sets-functions-relations/size-of-sets]{comparing-size}
\olimport*[sets-functions-relations/size-of-sets]{schroder-bernstein}
\olimport*[history/set-theory]{cantor-plane}
\olimport*[history/set-theory]{hilbert-curve}

\OLEndChapterHook

\olsetchapter{nosection}
\chapter{算术化}

在\olref[his][set][]{chap}中，我们考察了为何需要集合论的一些历史背景。\Crefrange{sfr:set::chap}{sfr:siz::chap}则探讨了朴素集合论的一些原理。所以，我们现在从朴素的角度理解了如何构造关系、函数以及比较集合的大小，\emph{诸如此类的内容}。

有了这些基础，我们便可以探讨集合论的一些早期成果，即在某种意义上将数学的大部分内容归约到集合论和算术。这就是本章的目标。

\olimport*[sets-functions-relations/arithmetization]{arithmetization}

\chapter{无限集}

在上一章中，我们展示了如何基于朴素集合论和自然数来构造一系列东西——整数、有理数和实数。本章要探讨的是：能否用集合论来构造自然数集\emph{本身}？

\olimport*[sets-functions-relations/infinite]{infinite}
\olresetchapter

\OLEndPartHook

\olpart{sth}{迭代概念}

\section*{\Olref[sth][][]{part}导言}

\Olref[sfr][][]{part}以朴素、非形式化的方式讨论了集合。现在是时候将其严格化，并提供一个关于集合的形式理论了。这正是\olref[sth][][]{part}的目标。

我们的形式理论是一个一阶理论，只包含一个二元谓词 $\in$。我们将规定几条支配属于关系的行为的公理。然而，我们只会在需要时引入这些公理，并在遇到它们时思考如何为其辩护。因此，我们将在接下来的各章中\emph{逐步}引入这些公理。

不过，让读者在一处看到所有公理的列表或许会有所帮助。所以，以下是我们在\olref[sth][][]{part}中将考虑的\emph{所有}公理。如同在\olref[sfr][][]{part}中一样，我们选择大小写字母仅出于可读性的考虑：

\begin{defish}
\emph{外延性公理。}
\[\forall A \forall B(\forall x(x \in A \liff x \in B) \lif A = B)\]
\end{defish}

\begin{defish}
\emph{并集公理。}
\[
\forall A \exists U \forall x(x \in U \liff (\exists b \in A)x \in b),
\]
即，对任何集合~$A$，$\bigcup A$ 存在。
\end{defish}

\begin{defish}
\emph{对集公理。}
\[
\forall a \forall b \exists P \forall x (x \in P \liff
(x = a \lor x = b)),
\]
即，对任何 $a$ 和~$b$，$\{a, b\}$ 存在。
\end{defish}

\begin{defish}
\emph{幂集公理。}
\[
\forall A \exists P \forall x(x \in P \liff (\forall z \in x)z \in A),
\]
即，对任何集合~$A$，$\Pow{A}$ 存在。
\end{defish}

\begin{defish}
\emph{无穷公理。}
\begin{multline*}
\exists I((\exists o \in I)\forall x(x \notin o) \land {}\\
(\forall x \in I)(\exists s \in I)\forall z(z \in s \liff (z \in x \lor z = x))),
\end{multline*}
即，存在一个集合，它以 $\emptyset$ 为元素且在 $x \mapsto x \cup \{x\}$ 下封闭。
\end{defish}

\begin{defish}
\emph{基础公理。}
\[
\forall A(\forall x\, x \notin A \lor (\exists b \in A)(\forall x \in A)x \notin b),
\]
即，$A = \emptyset$ 或 $(\exists b \in A)A\cap b = \emptyset$。
\end{defish}

\begin{defish}
\emph{良序公理。}
对每个集合 $A$，存在一个关系能良序 $A$。
（将其用一阶逻辑写出过于繁琐，不值得费心。）
\end{defish}

\begin{defish}
\emph{分离公理模式。}
对任何不包含``$S$''的公式 $\phi(x)$：
\[
\forall A \exists S \forall x(x \in S \liff (\phi(x) \land x \in A)),
\]
即，对任何集合~$A$，$\Setabs{x \in A}{\phi(x)}$ 存在。
\end{defish}

\begin{defish}
\emph{替换公理模式。}
对任何不包含``$B$''的公式 $\phi(x, y)$：
\[
\forall A((\forall x \in A)\lexists![y][\phi(x,y)] \lif \exists B \forall y (y \in B \liff (\exists x \in A)\phi(x,y))),
\]
即，对任何 $A$，如果 $\phi$ ``具有函数性质''，则 $\Setabs{y}{(\exists x \in A)\phi(x,y)}$ 存在。
\end{defish}

在两个模式中，公式都可能包含参数。事实上，在整个\olref[sth][][]{part}中，我们都遵循一个惯例：任何公式都可以包含参数。（关于公式可能包含参数的含义，请参见\olref[sfr][infinite][induction]{sec}以作复习。）

\begin{description}
\item[$\Zminus$] 包括外延性、并集、对集、幂集、无穷公理和分离公理模式。
\item[$\Z$] 指 $\Zminus$ 加上基础公理。
\item[$\ZFminus$] 指 $\Z$ 加上替换公理模式。
\item[$\ZF$] 指 $\ZFminus$ 加上基础公理。
\item[$\ZFC$] 指 $\ZF$ 加上良序公理。
\end{description}

\olimport*[set-theory/story]{story}

\chapter{迈向 $\Z$ 的步骤}

在前一章中，我们考察了集合的迭代概念。在本章中，我们将尝试从中提取出 Zermelo 集合论（即~$\Z$）的大部分公理。这一方法\emph{完全}受到了 \citet{Boolos1971}、\citet{Scott1974} 和 \citet{Shoenfield:AST} 的启发。

\olsetchapter{nosection}
\olimport*[set-theory/z]{z}
\olresetchapter

\olimport*[set-theory/ordinals]{ordinals}

\olimport*[set-theory/spine]{spine}

\olimport*[set-theory/replacement]{replacement}

\olimport*[set-theory/ord-arithmetic]{ord-arithmetic}

\olimport*[set-theory/cardinals]{cardinals}

\chapter{基数算术}

在 \olref[sth][cardinals][]{chap} 章中，我们发展了一套基数理论。下一步是概述基数算术理论。本章简要总结一些基本事实，然后指出其中的一些困难，而这些困难本身既引人入胜，又富有哲学意蕴。

\olsetchapter{nosection}
\olimport*[set-theory/card-arithmetic]{card-arithmetic}
\olresetchapter

\olimport*[set-theory/choice]{choice}

\stopproblems
\def\ifproblems#1{}

\appendix

\def\figurename{图}
\setlength{\olphotowidth}{.45\textwidth}

\chapter{传记}

\olimport*[history/biographies]{georg-cantor}

\olimport*[history/biographies]{kurt-goedel}

\olimport*[history/biographies]{bertrand-russell}

\olimport*[history/biographies]{alfred-tarski}

\olimport*[history/biographies]{ernst-zermelo}

\backmatter

\clearpage

\photocredits

\bibliographystyle{\olpath/bib/natbib-oup}
\bibliography{\olpath/bib/open-logic}

\olimport*{\olpath/content/open-logic-about}